\renewcommand{\theposttitle}{What most people get wrong about iteration and change}
\renewcommand{\thepostdate}{2024-02-13}
\vspace*{\fill}
\begin{center}
\begin{minipage}{0.92\textwidth}
\small
I’ve spoken before about the importance of iteration to bridge the gap between your ability and taste.

Today, I want to illustrate (🥁) \textbf{a common mistake} people make when they iterate: \textbf{focusing too heavily on “improving” the image}.

Of course we want the image to be better, but starting with that mindset actually blocks the creative process.

I tell my students all the time that \textbf{brainstorming is idea generation, not decision making. }

Starting with improvement in mind is decision making before the ideas have generated.

❌ how can I improve it? (explores only ideas you assume are good)\\
✅ how can articulate this idea in another way? (lets in all ideas, good and bad—we want ALL ideas to flow in our stream of consciousness)

\textbf{Improving is more a byproduct of trying out different articulations.}

Take this recent example of an iteration:
\end{minipage}
\end{center}
\vspace{4mm}
\begin{center}
\includegraphics[width=\textwidth, height=0.24\textheight, keepaspectratio]{C:/Users/James/Documents/GitHub/VisualIdeasBook/data/images/480892248086a40f0e4ee5610b5bcee0.png}
\end{center}
\vspace{2mm}
\begin{center}
\begin{minipage}{0.92\textwidth}
\small
Made in August 2022
\end{minipage}
\end{center}
\vspace{5mm}
\begin{center}
\includegraphics[width=\textwidth, height=0.24\textheight, keepaspectratio]{C:/Users/James/Documents/GitHub/VisualIdeasBook/data/images/8fd1a2ed542c39db5205c6430fd692ba.png}
\end{center}
\vspace{2mm}
\begin{center}
\begin{minipage}{0.92\textwidth}
\small
The nuance of these images are completely different but the core message is the same: address problems as they form to avoid it getting so convoluted that they become difficult to manage.

For me, one of the best parts of thinking in visual metaphors is the forcing function of having to get crystal clear on what you’re trying to say in your content.

\textbf{CLARITY = iteration x time}

Ultimately, whether version 2 is better or worse is irrelevant. The act of creating a different articulation opens another door to how deeply you think about the idea.

When you’re ready and if you’re curious to learn how to make visual metaphors like I do, \textbf{sign up for the waitlist for the next Thinking in Visual Metaphors cohort. }
\end{minipage}
\end{center}
\vspace*{\fill}